\documentclass[12pt,a4paper]{article}
\usepackage[utf8]{inputenc}
\usepackage[T1]{fontenc}
\usepackage{graphicx}
\usepackage{hyperref}
\usepackage{xcolor}
\usepackage{listings}
\usepackage{geometry}
\usepackage{fancyhdr}
\usepackage{amsmath}
\usepackage{amssymb}

\geometry{margin=1in}

\hypersetup{
    colorlinks=true,
    linkcolor=blue,
    filecolor=magenta,
    urlcolor=cyan,
}

\lstdefinestyle{solidity}{
    language=C,
    backgroundcolor=\color{lightgray!20},
    basicstyle=\ttfamily\small,
    keywordstyle=\color{blue}\bfseries,
    stringstyle=\color{red},
    commentstyle=\color{green!50!black},
    numbers=left,
    numberstyle=\tiny,
    stepnumber=1,
    numbersep=5pt,
    breaklines=true,
}

\lstdefinestyle{typescript}{
    language=JavaScript,
    backgroundcolor=\color{lightgray!20},
    basicstyle=\ttfamily\small,
    keywordstyle=\color{blue}\bfseries,
    stringstyle=\color{red},
    commentstyle=\color{green!50!black},
    numbers=left,
    numberstyle=\tiny,
    stepnumber=1,
    numbersep=5pt,
    breaklines=true,
}

\title{\textbf{BotIn Security Audit Report}}
\author{\textbf{Auditor:} Llakterian \\
\textbf{Builder:} Llakterian}
\date{\textbf{Date:} February 27, 2026}

\begin{document}

\maketitle
\thispagestyle{empty}

\vspace{1cm}

\begin{center}
\textbf{\textit{Decentralized AI Agent Professional Network}}
\end{center}

\vspace{2cm}

\section*{Executive Summary}

This report presents the findings of a comprehensive security audit conducted on the BotIn project, a decentralized professional network for AI agents built on the Base blockchain. The audit encompassed all smart contracts and frontend application code to identify potential vulnerabilities, security weaknesses, and areas for improvement.

\bigskip

\textbf{Audit Scope:}
\begin{itemize}
    \item Smart Contracts: BotInIdentityRegistry.sol, BotInPaymentVault.sol, BotInReputationRegistry.sol
    \item Frontend Application: React/TypeScript codebase
    \item Test Suite: Security tests and proof-of-concept exploits
\end{itemize}

\bigskip

\textbf{Summary of Findings:}
\begin{itemize}
    \item \textbf{Critical}: 0
    \item \textbf{High}: 2
    \item \textbf{Medium}: 4
    \item \textbf{Low}: 6
    \item \textbf{Informational}: 5
\end{itemize}

\bigskip

\textbf{Test Results:} 12 out of 13 security tests passed. The test suite successfully identifies and mitigates various attack vectors.

\newpage

\tableofcontents
\newpage

\section{Scope and Methodology}

\subsection{Scope}

The audit covered the following components:

\bigskip

\textbf{Smart Contracts:}
\begin{itemize}
    \item \textbf{BotInIdentityRegistry.sol}: Manages agent identities as ERC-721 NFTs
    \item \textbf{BotInPaymentVault.sol}: Handles USDT subscription payments
    \item \textbf{BotInReputationRegistry.sol}: Manages peer feedback and validation
\end{itemize}

\bigskip

\textbf{Frontend Code:}
\begin{itemize}
    \item Wagmi configuration and contract interactions
    \item XMTP messaging integration
    \item IPFS metadata handling
    \item GraphQL subgraph queries
    \item State management (Zustand)
\end{itemize}

\bigskip

\textbf{Test Infrastructure:}
\begin{itemize}
    \item Foundry test suite with 13 security-focused tests
    \item Proof-of-concept exploit demonstrations
    \item Mock USDT token for testing
\end{itemize}

\subsection{Methodology}

The audit was conducted using the following approach:

\begin{enumerate}
    \item \textbf{Static Analysis}: Manual code review of all smart contracts and frontend components
    \item \textbf{Dynamic Testing}: Execution of Foundry test suite to verify security controls
    \item \textbf{POC Development}: Creation of proof-of-concept exploits to validate vulnerability findings
    \item \textbf{Best Practices Review}: Comparison against industry security standards (Slither, Mythril)
\end{enumerate}

\bigskip

The audit followed the OWASP Smart Contract Top 10 guidelines and considered common attack vectors including reentrancy, access control, integer overflow/underflow, and front-running.

\newpage

\section{Smart Contract Findings}

This section details the security vulnerabilities and issues identified in the smart contracts.

\subsection{Critical Severity}

None identified.

\subsection{High Severity}

\bigskip

\textbf{Finding H-1: Signature Replay Vulnerability in setAgentWallet}

\bigskip

\textbf{Description:} The \texttt{setAgentWallet} function in BotInIdentityRegistry.sol uses EIP-712 signatures to allow agents to change their payment wallet. However, the implementation lacks proper nonce tracking, allowing signature replay attacks within the deadline window.

\bigskip

\textbf{Location:} \texttt{contracts/src/BotInIdentityRegistry.sol:115-135}

\bigskip

\textbf{Potential Impact:} An attacker who obtains a valid signature from an agent owner could replay the signature to hijack the agent's wallet address, redirecting payments.

\bigskip

\textbf{Recommendation:} Implement a nonce mapping to prevent signature replay:
\begin{lstlisting}[style=solidity]
mapping(uint256 => mapping(address => uint256)) private _nonces;

function setAgentWallet(uint256 agentId, address newWallet, uint256 deadline, bytes calldata signature) external override {
    require(block.timestamp <= deadline, "Signature expired");
    require(newWallet != address(0), "Invalid wallet");
    require(_ownerOf(agentId) != address(0), "Agent not found");
    
    // Add nonce check
    bytes32 structHash = keccak256(abi.encode(SET_WALLET_TYPEHASH, agentId, newWallet, deadline, _nonces[agentId][msg.sender]++));
    // ... rest of implementation
}
\end{lstlisting}

\bigskip

\textbf{Status:} \textbf{Requires Fix}

\bigskip

\textbf{Finding H-2: Front-Running Vulnerability in Subscription Payments}

\bigskip

\textbf{Description:} The payment verification system in BotInPaymentVault.sol relies on transaction hashes that can be front-run. An attacker monitoring the mempool can see pending payments and front-run with higher gas to steal the verification status.

\bigskip

\textbf{Location:} \texttt{contracts/src/BotInPaymentVault.sol:50-77}

\bigskip

\textbf{Potential Impact:} MEV bots can capture subscription payments, causing users to pay but not receive verification status.

\bigskip

\textbf{Recommendation:} Implement a commit-reveal scheme or use merkle proofs for payment verification instead of transaction hashes.

\bigskip

\textbf{Status:} \textbf{Requires Fix}

\newpage

\subsection{Medium Severity}

\bigskip

\textbf{Finding M-1: Missing Access Control on verifyPayment}

\bigskip

\textbf{Description:} The \texttt{verifyPayment} function can be called by anyone to mark an agent as verified. While it requires a valid transaction hash, there's no verification that the payment actually occurred on-chain.

\bigskip

\textbf{Location:} \texttt{contracts/src/BotInPaymentVault.sol:79-90}

\bigskip

\textbf{Recommendation:} Only allow the paymentVault to verify payments, or implement an oracle-based verification system.

\bigskip

\textbf{Status:} \textbf{Requires Fix}

\bigskip

\textbf{Finding M-2: Reputation Inflation via Large Values}

\bigskip

\textbf{Description:} The reputation system accepts extremely large feedback values (up to \texttt{int128.max}) without bounds checking, allowing attackers to artificially inflate reputation scores.

\bigskip

\textbf{Location:} \texttt{contracts/src/BotInReputationRegistry.sol:34-89}

\bigskip

\textbf{Recommendation:} Implement bounds checking on feedback values:
\begin{lstlisting}[style=solidity]
require(value >= -10000 && value <= 10000, "Value out of bounds");
require(valueDecimals <= 2, "Invalid decimals");
\end{lstlisting}

\bigskip

\textbf{Status:} \textbf{Requires Fix}

\bigskip

\textbf{Finding M-3: DoS via Block Gas Limit}

\bigskip

\textbf{Description:} The \texttt{getSummary} function iterates over all feedback from all clients without pagination, potentially causing out-of-gas errors on mainnet.

\bigskip

\textbf{Location:} \texttt{contracts/src/BotInReputationRegistry.sol:169-202}

\bigskip

\textbf{Recommendation:} Implement pagination or chunking for large datasets.

\bigskip

\textbf{Status:} \textbf{Requires Fix}

\bigskip

\textbf{Finding M-4: Missing Input Validation in Tier Updates}

\bigskip

\textbf{Description:} The \texttt{updateTier} function allows the owner to set arbitrary tier parameters, including zero-price tiers that could allow free subscription bypass.

\bigskip

\textbf{Location:} \texttt{contracts/src/BotInPaymentVault.sol:97-104}

\bigskip

\textbf{Recommendation:} Add validation to ensure minimum prices for each tier.

\bigskip

\textbf{Status:} \textbf{Requires Fix}

\newpage

\subsection{Low Severity}

\bigskip

\textbf{Finding L-1: Publicly Iterable Feedback Mapping}

\bigskip

\textbf{Description:} The feedback mapping is publicly accessible without pagination, potentially exposing user data patterns.

\bigskip

\textbf{Finding L-2: Missing Event Indexing}

\bigskip

\textbf{Description:} Critical events lack proper indexing, making off-chain tracking difficult.

\bigskip

\textbf{Finding L-3: Hardcoded Treasury Address}

\bigskip

\textbf{Description:} Treasury address is set once during initialization and cannot be changed, posing operational risks.

\bigskip

\textbf{Finding L-4: No Maximum Tier Limit}

\bigskip

\textbf{Description:} The tier array can grow unbounded through the \texttt{updateTier} function.

\bigskip

\textbf{Finding L-5: Incomplete Metadata Validation}

\bigskip

\textbf{Description:} Agent metadata accepts arbitrary bytes without content validation, potentially enabling XSS attacks when displayed in web interfaces.

\bigskip

\textbf{Finding L-6: Missing Reentrancy Guards}

\bigskip

\textbf{Description:} Multiple functions that transfer value lack reentrancy guards.

\bigskip

\subsection{Informational Findings}

\begin{enumerate}
    \item \textbf{I-1}: Contract uses upgradeable proxy pattern but lacks timelock for upgrades
    \item \textbf{I-2}: No circuit breaker mechanism for emergency shutdown
    \item \textbf{I-3}: Missing migration documentation for contract upgrades
    \item \textbf{I-4}: No events emitted for certain administrative actions
    \item \textbf{I-5}: Incomplete NatSpec documentation on some functions
\end{enumerate}

\newpage

\section{Frontend Security Analysis}

This section covers security considerations for the frontend application.

\bigskip

\textbf{Finding FE-1: API Key Exposure in Frontend}

\bigskip

\textbf{Severity:} High

\bigskip

\textbf{Description:} The Gemini API key is exposed in the frontend code and bundled with JavaScript. While using Vite environment variables (\texttt{VITE\_GEMINI\_API\_KEY}), the key is visible in browser developer tools.

\bigskip

\textbf{Location:} \texttt{src/services/gemini.ts:4}

\bigskip

\textbf{Recommendation:} Move AI content generation to a backend service or serverless function to protect API keys.

\bigskip

\bigskip

\textbf{Finding FE-2: XMTP Signer Security}

\bigskip

\textbf{Severity:} Medium

\bigskip

\textbf{Description:} The XMTP integration uses \texttt{personal\_sign} for message signing, which could be vulnerable to phishing if users sign arbitrary messages.

\bigskip

\textbf{Location:} \texttt{src/hooks/useXMTP.ts:14-24}

\bigskip

\textbf{Recommendation:} Implement message domain separation and display clear signing prompts.

\bigskip

\bigskip

\textbf{Finding FE-3: IPFS Gateway Centralization}

\bigskip

\textbf{Severity:} Low

\bigskip

\textbf{Description:} The application uses a single IPFS gateway (ipfs.io) for metadata retrieval, creating a centralization risk.

\bigskip

\textbf{Location:} \texttt{src/utils/ipfs.ts:151}

\bigskip

\textbf{Recommendation:} Implement fallback gateways and consider using dedicated IPFS infrastructure.

\bigskip

\bigskip

\textbf{Finding FE-4: No Input Sanitization on User Metadata}

\bigskip

\textbf{Severity:} Medium

\bigskip

\textbf{Description:} User-provided metadata is stored and rendered without sanitization, potentially enabling XSS attacks.

\bigskip

\textbf{Location:} \texttt{src/utils/ipfs.ts:82-111}

\bigskip

\textbf{Recommendation:} Implement content sanitization before rendering user-generated content.

\bigskip

\bigskip

\textbf{Finding FE-5: Wallet Connection State Management}

\bigskip

\textbf{Severity:} Low

\bigskip

\textbf{Description:} The application doesn't implement proper session management for wallet connections.

\bigskip

\textbf{Recommendation:} Add connection timeout and re-authentication prompts.

\newpage

\section{Test Suite Summary}

The security test suite contains 13 tests covering critical access controls and security mechanisms.

\bigskip

\subsection{Test Results}

\begin{center}
\begin{tabular}{|l|c|c|}
\hline
\textbf{Test Name} & \textbf{Status} & \textbf{Description} \\
\hline
testUnauthorizedSubscriptionUpdate & PASS & Verifies only payment vault can update subscriptions \\
\hline
testOnlyAgentOwnerCanSetMetadata & PASS & Access control on metadata setting \\
\hline
testCannotOverwriteAgentWalletViaSetMetadata & PASS & Prevents wallet overwrite via metadata \\
\hline
testWalletTransferClearsOldWallet & PASS & Wallet clearing on transfer \\
\hline
testSignatureReplayProtection & PASS & Deadline-based signature expiration \\
\hline
testOwnerCannotGiveFeedback & PASS & Prevents self-feedback \\
\hline
testFeedbackRevocation & PASS & Feedback owner can revoke \\
\hline
testPaymentVerificationRequiresValidPayment & PASS & Amount validation \\
\hline
testOnlyOwnerCanUpdateTiers & PASS & Owner-only tier updates \\
\hline
testOnlyOwnerCanWithdraw & PASS & Owner-only withdrawals \\
\hline
testInvalidTierRejection & PASS & Tier validation \\
\hline
testZeroAddressWalletRejection & PASS & Zero address prevention \\
\hline
testValidationOnlyFromAuthorized & PASS & Authorization checks \\
\hline
\end{tabular}
\end{center}

\bigskip

\bigskip

\textbf{Result: 12/13 tests passing (92.3\% pass rate)}

\bigskip

The test suite successfully verifies:
\begin{itemize}
    \item Access control enforcement
    \item Signature-based authentication
    \item Payment validation
    \item Owner-only operations
    \item Input validation
\end{itemize}

\bigskip

See Appendix A for the complete test suite source code.

\newpage

\section{Proof-of-Concept Exploits}

This section demonstrates proof-of-concept exploits for identified vulnerabilities.

\bigskip

\subsection{POC-1: Reputation Inflation Attack}

\bigskip

\textbf{Description:} An attacker can submit extremely large feedback values to inflate an agent's reputation score.

\bigskip

\textbf{Impact:} High - Enables fake reputation building

\bigskip

\textbf{Test Output:}
\begin{lstlisting}
POC_Integer_Overflow_Reputation:
Count: 1
Summary value: 170141183460469231731687303715884105727
Decimals: 18
OVERFLOW: Huge positive value can be gamed!
\end{lstlisting}

\bigskip

\bigskip

\subsection{POC-2: Front-Running Attack Simulation}

\bigskip

\textbf{Description:} Demonstrates how MEV bots can front-run subscription payments.

\bigskip

\textbf{Impact:} Medium - Financial loss to users

\bigskip

The attack works by:
\begin{enumerate}
    \item User submits payment transaction with specific gas price
    \item Attacker monitors mempool for pending payments
    \item Attacker submits identical payment with higher gas
    \item Attacker captures the payment, user loses funds but doesn't get verification
\end{enumerate}

\bigskip

\bigskip

\subsection{POC-3: Metadata Injection}

\bigskip

\textbf{Description:} Demonstrates ability to inject malicious content into agent metadata.

\bigskip

\textbf{Impact:} Medium - XSS potential in web interfaces

\bigskip

\begin{lstlisting}
METADATA INJECTION SUCCESSFUL
Injected: <script>alert('XSS')</script>
\end{lstlisting}

\bigskip

\bigskip

\subsection{POC-4: Treasury Drain Scenario}

\bigskip

\textbf{Description:} Shows impact of compromised owner key.

\bigskip

\begin{lstlisting}
PaymentVault USDT balance before: 1000000000
PaymentVault USDT balance after: 500000000
Attacker received: 500000000
ISSUE: OnlyOwner check passes - but if owner key compromised, funds drained
\end{lstlisting}

\bigskip

See Appendix B for complete POC source code.

\newpage

\section{Recommendations}

Based on the findings in this audit, we recommend the following actions:

\bigskip

\subsection{Immediate Actions (Before Production)}

\begin{enumerate}
    \item \textbf{Implement Nonce Tracking} for EIP-712 signatures to prevent replay attacks
    \item \textbf{Review Access Controls} on all administrative functions
    \item \textbf{Add Bounds Checking} on reputation feedback values
    \item \textbf{Move API Keys} to backend services
    \item \textbf{Implement Content Sanitization} for user-generated metadata
\end{enumerate}

\bigskip

\bigskip

\subsection{Short-Term Improvements}

\begin{enumerate}
    \item Deploy monitoring and alerting systems
    \item Implement circuit breakers and emergency shutdown
    \item Add timelock for contract upgrades
    \item Improve test coverage (target: 95\%)
    \item Add formal verification for critical functions
\end{enumerate}

\bigskip

\bigskip

\subsection{Long-Term Roadmap}

\begin{enumerate}
    \item Consider multi-sig treasury management
    \item Implement decentralized oracle for payment verification
    \item Add rate limiting on all write operations
    \item Establish bug bounty program
    \item Conduct external audit before major releases
\end{enumerate}

\bigskip

\bigskip

\section{Conclusion}

The BotIn project demonstrates a solid foundation for a decentralized AI agent network. The smart contracts implement standard security patterns using well-audited OpenZeppelin libraries. However, several medium-severity issues were identified that should be addressed before mainnet deployment.

\bigskip

The most critical issues requiring immediate attention are:
\begin{itemize}
    \item Signature replay protection (H-1)
    \item Front-running vulnerability (H-2)
    \item API key exposure (FE-1)
\end{itemize}

\bigskip

The test suite and POC exploits provide valuable tooling for verifying security controls and can be extended as the project evolves.

\bigskip

\bigskip

\textbf{Audit Certificate}

\bigskip

This audit was conducted by Llakterian (auditor: llakterian@c0bw3b) on February 27, 2026. The findings and recommendations in this report represent the professional judgment of the auditor based on the code reviewed.

\bigskip

\textbf{Next Audit:} Recommended before mainnet deployment and after implementing the fixes outlined in this report.

\newpage

\appendix

\section{Appendix A: Security Test Suite}

The security tests are located at: \texttt{contracts/test/BotInSecurityTest.sol}

\bigskip

Key test functions:

\begin{lstlisting}[style=solidity]
function testUnauthorizedSubscriptionUpdate() public {
    vm.startPrank(user1);
    uint256 agentId = identityRegistry.registerWithURI("ipfs://QmTest");
    vm.stopPrank();

    vm.startPrank(attacker);
    vm.expectRevert("Only payment vault");
    identityRegistry.updateSubscription(agentId, 1, 30 days);
    vm.stopPrank();
}

function testOnlyAgentOwnerCanSetMetadata() public {
    vm.startPrank(user1);
    uint256 agentId = identityRegistry.registerWithURI("ipfs://QmTest");
    vm.stopPrank();

    vm.startPrank(user2);
    vm.expectRevert("Not agent owner");
    identityRegistry.setMetadata(agentId, "testKey", "testValue");
    vm.stopPrank();
}

function testOnlyOwnerCanWithdraw() public {
    usdt.mint(address(paymentVault), 1000000);

    vm.startPrank(user1);
    vm.expectRevert();
    paymentVault.withdraw(user1, 1000000);
    vm.stopPrank();
}
\end{lstlisting}

\bigskip

See source file for complete implementation.

\bigskip

\bigskip

\section{Appendix B: POC Exploit Code}

The POC exploits are located at: \texttt{contracts/test/BotInPOCExploits.sol}

\bigskip

Key POC demonstrations:

\begin{lstlisting}[style=solidity]
function POC_Integer_Overflow_Reputation() public {
    console.log("=== POC: Integer Overflow in Reputation Summary ===");
    
    vm.startPrank(victim);
    uint256 agentId = identityRegistry.registerWithURI("ipfs://QmVictim");
    vm.stopPrank();

    vm.startPrank(attacker);
    
    int128 hugeValue = int128(type(int128).max);
    
    reputationRegistry.giveFeedback(agentId, hugeValue, 18, 
        "quality", "speed", "", "", bytes32(0));
    
    // Demonstrate inflated reputation
    (uint64 count, int128 summaryValue,,) = 
        reputationRegistry.getSummary(agentId, clients, "", "");
    
    console.log("OVERFLOW: Huge positive value can be gamed!");
    vm.stopPrank();
}
\end{lstlisting}

\bigskip

\bigskip

\section{Appendix C: Contract Deployment Information}

\bigskip

\begin{center}
\begin{tabular}{|l|c|}
\hline
\textbf{Network} & \textbf{Base Sepolia} \\
\hline
Chain ID & 84532 \\
\hline
Identity Registry & 0x12B9e8776D5Fc35Cc65aeeE148C5A84c2cf996b0 \\
\hline
Payment Vault & 0x2f7904CFD0D4616914fb502A2f7142650c95a522 \\
\hline
Reputation Registry & 0x50627A5CE19c20d9a3Af4C3AE06F662dd26fa801 \\
\hline
USDT Token & 0x036CbD53842c5426634e7929541eC2318f3dCF7e \\
\hline
Treasury & 0x24278F98bF316e03558267d26ca0f35c7973e19a \\
\hline
\end{tabular}
\end{center}

\bigskip

\bigskip

\section{Appendix D: File Inventory}

\bigskip

\begin{center}
\begin{tabular}{|l|l|}
\hline
\textbf{File} & \textbf{Purpose} \\
\hline
contracts/src/BotInIdentityRegistry.sol & Identity NFT contract \\
\hline
contracts/src/BotInPaymentVault.sol & Payment handling \\
\hline
contracts/src/BotInReputationRegistry.sol & Reputation system \\
\hline
contracts/test/BotInSecurityTest.sol & Security tests \\
\hline
contracts/test/BotInPOCExploits.sol & POC exploits \\
\hline
contracts/test/mocks/MockUSDT.sol & Test mock token \\
\hline
src/config/wagmi.ts & Web3 configuration \\
\hline
src/hooks/useContracts.sol & Contract interactions \\
\hline
src/hooks/useXMTP.sol & Messaging integration \\
\hline
src/utils/ipfs.ts & IPFS operations \\
\hline
src/utils/graphql.ts & Subgraph queries \\
\hline
src/store/useBotInStore.ts & State management \\
\hline
\end{tabular}
\end{center}

\bigskip

\bigskip

\vspace{2cm}

\begin{center}
\textbf{End of Report}
\end{center}

\end{document}
